\chapter{Variables aléatoires}

\section{Variables aléatoires et leurs lois}

Soient $(E, \calE)$ et $(F, \calF)$ deux espaces probabilisables
($E, F$~: ensembles, $\calE, \calF$~: tribus sur $E, F$).

\begin{deff}
    Une application $f: E \to F$ est \textbf{mesurable} si pour tout $B \in \calF$, on a $f^{-1}(B) \in \calE$.
\end{deff}

Si $(\Omega, \calA, \bbP)$ est un espace de probabilité,
alors une variable aléatoire à valeurs dans $(E, \calE)$ est une application $X: (\Omega, \calA) \to (E, \calE)$.
Autrement dit, une variable aléatoire est simplement une fonction mesurable définie sur un espace de probabilité.

\begin{lem}
    Soient $(E, \calE), (F, \calF), (G, \calG)$ trois espaces probabilisables.
    Si $f: E \to F$ et $g: F \to G$ sont deux fonctions mesurables,
    alors $g \circ f: E \to G$ est aussi mesurable.
\end{lem}

\begin{proof}
    Si $f$ et $g$ sont mesurables,
    alors $\forall C \in \calG, g^{-1}(C) \in \calF, \forall B \in \calF, f^{-1}(B) \in \calE$.
    Alors $(g \circ f)^{-1}(C) = \{ x \in E \mid g(f(x)) \in C \}
        = \{ x \in E \mid f(x) \in g^{-1}(C) \} = f^{-1}(g^{-1}(C))$.
    $B := g^{-1}(C) \in \calF$, donc $f^{-1}(B) \in \calE$.
    Donc $\forall C \in \calG, (g \circ f)^{-1}(C) \in \calE$.
    Càd $g \circ f$ est mesurable.
\end{proof}

\begin{cor}
    Si $X: (\Omega, \calA, \bbP) \to (E, \calE)$ est une variable aléatoire,
    et $f: (E, \calE) \to (F, \calF)$ est une fonction mesurable,
    alors $f(X) := f \circ X$ est une variable aléatoire à valeurs dans $(F, \calF)$.
\end{cor}

\begin{rqq}
    Si $X: (\Omega, \calA, \bbP) \to (E, \calE)$ est une variable aléatoire,
    alors pour tout ensemble mesurable $B \in \calE$
    (``question mesurable par rapport à la valeur de la variable aléatoire''),
    l'ensemble $\{ X \in B \} := \{ \omega \in \Omega \mid X(\omega) \in B \}$ est un événement
    (i.e. mesurable dans $\calA$).
\end{rqq}

\begin{rqq}
    Rappel~: la tribu borélienne de $\bbR$ est définie par
    \[ B(\bbR) := \sigma \left(\{ ]-\infty, t] \mid t \in \bbR \}\right). \]
    Intuitivement, la tribu $B(\bbR)$ peut ``répondre'' à toutes les questions que vous pouvez construire à partir de
    ``est-ce que $x \leqslant t$''.
\end{rqq}

\begin{deff}
    Lorsque $X$ est une variable aléatoire à valeurs dans $(\bbR, B(\bbR))$,
    on l'appelle une variable aléatoire \textbf{réelle}.
\end{deff}

\begin{prop}
    Soit une application $X: (\Omega, \calA, \bbP) \to (\bbR, B(\bbR))$,
    alors $X$ est mesurable (càd $X$ est une variable aléatoire réelle) ssi
    \[ \{ X \leqslant t \} := \{ \omega \in \Omega \mid X(\omega) \leqslant t \} \in \calA. \]
\end{prop}

\begin{proof}
    ``$\Rightarrow$''~: Si $X$ est mesurable, alors
    \[ \{ X \leqslant t \} = X^{-1}(]-\infty, t]) \in \calA. \]

    ``$\Leftarrow$''~: Supposons $\{ X \leqslant t \} \in \calA, \forall t \in \bbR$.
    Soit $\calC = \{ B \subseteq \bbR \mid X^{-1}(B) \in \calA \}$.
    Alors~:
    \begin{itemize}
        \item $\calA$ est une tribu sur $\Omega$, donc $\Omega \in \calA
                  \implies \bbR \in \calC \implies X^{-1}(\bbR) = \Omega \in \calA$.
        \item Soit $(B_n)$ une suite d'événements dans $\calC$,
              donc pour tout $n$, $X^{-1}(B_n) \in \calA$.
              Donc $\bigcap\limits_{n \in \bbN} X^{-1}(B_n) = X^{-1} \left(\bigcap\limits_{n \in \bbN} B_n\right) \in \calA$,
              $\left(X^{-1}(B_n)\right)^C = X^{-1}(B_n^C) \in \calA$.
              Donc $\bigcap\limits_{n \in \bbN} B_n \in \calC$, $B_n^C \in \calC$.
              Donc $\calC$ est stable par intersection dénombrable et par complémentaire.
              $\calC$ est une tribu sur $\bbR$.
              Mais puisque $\{ X \leqslant t \} = X^{-1}(]-\infty, t]) \in \calA$,
              on a $\{ ]-\infty, t] \mid t \in \bbR \} \subseteq \calC$.
              Donc $\sigma(\{ ]-\infty, t] \mid t \in \bbR \}) \subseteq \calC$.
              Càd $B(\bbR) \subseteq \calC \iff \forall B \in B(\bbR), X^{-1}(B) \in \calA$
              $\iff X$ est une variable aléatoire réelle.
    \end{itemize}
\end{proof}

\begin{exem}
    $\ $
    \begin{enumerate}
        \item Toute application constante $X:
                  \begin{array}[t]{ccc}
                      \Omega & \to & E                                         \\
                      \omega & \tp & c \in E (\text{ne dépend pas de } \omega)
                  \end{array}
              $ est une variable aléatoire.

              Une telle variable aléatoire est dite ``déterministe''.
        \item Si $A \in \calA$ est un événement dans l'espace de probabilité $(\Omega, \calA, \bbP)$,
              alors $\bbI_A:
                  \begin{array}[t]{ccc}
                      \Omega & \to & \bbR                          \\
                      \omega & \tp & 1 \text{ si } \omega \in A    \\
                      \omega & \tp & 0 \text{ si } \omega \notin A
                  \end{array}
              $ est une variable aléatoire réelle.
        \item Si $([0, 1], B([0, 1]), \bbP)$ est un espace de probabilité sur $[0, 1]$,
              alors $\id_{[0, 1]}: ([0, 1], B([0,1])) \to ([0, 1], B([0, 1]))$ est une variable aléatoire.
              Cette variable aléatoire peut prendre une infinité non dénombrable de valeurs (i.e. $[0, 1]$).
    \end{enumerate}
\end{exem}

\begin{proof}
    $\ $
    \begin{enumerate}
        \item Pour tout $B \subseteq E$, on a $X^{-1}(B) =
                  \begin{cases}
                      \varnothing, & \text{si } c \notin B \\
                      \Omega,      & \text{si } c \in B
                  \end{cases}
              $

              $\varnothing$ et $\Omega$ appartiennent toujours à une tribu sur $\Omega$.
              Donc $X$ est mesurable.
        \item Soit $t \in \bbR$,
              \[ \bbI_A^{-1}(]-\infty, t]) =
                  \begin{cases}
                      \Omega,      & \text{si } t \geqslant 1     \\
                      A^C,         & \text{si } 0 \leqslant t < 1 \\
                      \varnothing, & \text{si } t < 0
                  \end{cases}
                  \in \calA \] pour tout $t \in \bbR$.
              D'après la proposition précédente, $\bbI_A: (\Omega, \calA, \bbP) \to (\bbR, B(\bbR))$ est mesurable,
              donc c'est une variable aléatoire réelle.
    \end{enumerate}
\end{proof}

\begin{rqq}[Notation]
    Soit $X: (\Omega, \calA, \bbP) \to (E, \calE)$ une variable aléatoire à valeurs dans $(E, \calE)$.
    \begin{itemize}
        \item Pour $B \in \calE$, on écrit
              \[ \{ X \in B \} := \{ \omega \in \Omega \mid X(\omega) \in B \} = X^{-1}(B). \]
        \item Pour une variable aléatoire réelle (i.e. $(E, \calE) = (\bbR, B(\bbR))$), on écrit
              \begin{align*}
                  \{ X > x_0 \}           & := \{ \omega \in \Omega \mid X(\omega) > x_0 \},           \\
                  \{ a < X \leqslant b \} & := \{ \omega \in \Omega \mid a < X(\omega) \leqslant b \}.
              \end{align*}
    \end{itemize}
\end{rqq}

\begin{deff}
    Soit $X: (\Omega, \calA, \bbP) \to (E, \calE)$ une variable aléatoire.
    Sa loi, notée $\bbP_X$, est l'application
    \[ \bbP_X:
        \begin{array}[t]{ccc}
            \calE & \to & [0, 1]        \\
            B     & \tp & \bbP(X \in B)
        \end{array}
    \]

    ($\bbP_X$ est bien définie car pour tout $B \in \calE$, on a $\{ X \in B \} = X^{-1}(B) \in \calA$.
    Donc $\bbP(X^{-1}(B))$ est bien définie.)
\end{deff}

\begin{thm}
    La loi $\bbP_X$ est une mesure de probabilité sur $(E, \calE)$.
\end{thm}

\begin{proof}
    $\ $
    \begin{enumerate}
        \item $\bbP_X(E) = \bbP(X \in E) = \bbP(\Omega) = 1$.
        \item Soit $(E_n)_{n \in \bbN}$ une suite d'éléments de $\calE$ deux à deux disjoints.
              Alors $(X^{-1}(E_n))_{n \in \bbN}$ est une suite d'ensembles deux à deux disjoints dans $\calA$.
              Donc par la $\sigma$-additivité de $\bbP$,
              \[ \bbP \left(\bigcup\limits_{n \in \bbN} X^{-1}(E_n) \right) = \sum\limits_{n \in \bbN} \bbP(X^{-1}(E_n)). \]
              Donc $\bbP_X \left(\bigcup\limits_{n \in \bbN} E_n\right) = \sum\limits_{n \in \bbN} \bbP_X(E_n)$.
              Càd $\bbP_X$ est aussi $\sigma$-additive (sur $(E, \calE)$).
              Donc $\bbP_X$ est une probabilité sur $(E, \calE)$.
    \end{enumerate}
\end{proof}

\begin{rqq}
    Pour n'importe quelle mesure de probabilité $\mu$ sur l'espace probabilisable $(E, \calE)$,
    il existe une variable aléatoire $X: (\Omega, \calA, \bbP) \to (E, \calE)$ telle que $\mu = \bbP_X$.

    Il suffit de considérer $(\Omega, \calA, \bbP) = (E, \calE, \mu)$ et $X = \id_E$.
    i.e. $X$ est une variable aléatoire à valeurs dans $(E, \calE)$ et la loi de $X$ est $\mu$.
\end{rqq}

\begin{deff}
    Soient $X, Y$ deux variables aléatoires à valeurs dans le même espace probabilisable $(E, \calE)$.
    On dit que $X$ et $Y$ ont \textbf{la même loi} (ou~: $X$ et $Y$ sont \textbf{équiréparties}) si $\bbP_X = \bbP_Y$.
    Dans ce cas, on écrit $X \overset{\text{loi}}{=} Y$ (``$X$ est égale à $Y$'' en loi) ou
    $X \overset{\text{(d)}}{=} Y$ (``$X$ est égale à $Y$'' en distribution).
\end{deff}

\begin{rqq}
    $X$ et $Y$ ne sont pas nécessairement définies sur le même espace de probabilité $(\Omega, \calA, \bbP)$.
\end{rqq}

\subsection{Variable aléatoire discrète}

\begin{deff}
    Une variable aléatoire $X: \Omega \to E$ est \textbf{discrète},
    si son image $X(\Omega) = \{ X(\omega) \mid \omega \in \Omega \}$ est au plus dénombrable.

    Si $X(\Omega)$ est finie, $X$ est appelée une variable aléatoire finie.
\end{deff}

\begin{rqq}[Hypothèse]
    Dans la suite,
    on suppose que les singletons $\{ x \}$ ($x \in E$) dans l'espace d'arrivée $(E, \calE)$ sont mesurables.
    (càd $\forall x \in E, \{ x \} \in \calE$).
    En particulier, cette hypothèse est vraie si $E = \bbR^n$ et $\calE = B(\bbR^n)$ ou $\calE = \calP(E)$.
\end{rqq}

\begin{prop}
    Soit $X$ une variable aléatoire discrète à valeurs dans $(E, \calE)$.
    Alors $\forall B \in \calE, \bbP(X \in B) = \sum\limits_{x \in B \cap X(\Omega)} \bbP(X = x)$.
\end{prop}

\begin{proof}
    On a
    \[ \{ X \in B \} = \{ \omega \in \Omega \mid X(\omega) \in B \} =
        \bigcup_{x \in B \cap X(\Omega)} \{ \omega \in \Omega \mid X(\omega) = x \}. \]
    Les ensembles $\{ \omega \in \Omega \mid X(\omega) = x \}$ (pour $x \in B \cap X(\Omega)$) sont deux à deux disjoints,
    donc par la $\sigma$-additivité de $\bbP$,
    \[ \bbP(X \in B) = \sum_{x \in B \cap X(\Omega)} \bbP(X = x). \qedhere \]
\end{proof}

\begin{exem}
    Si $X = \bbI_A$ (avec $A \in \calA$),
    alors
    \begin{itemize}
        \item $\bbP(X = 1) = \bbP(A)$.
        \item $\bbP(X = 0) = \bbP(A^C) = 1 - \bbP(A)$.
        \item $X$ est une variable réelle et discrète.
    \end{itemize}
\end{exem}

\subsection{Fonction de répartition}

\begin{deff}
    Soit $X$ une variable aléatoire réelle.
    La fonction de répartition de $X$ est l'application
    \[ F_X:
        \begin{array}[t]{ccc}
            \bbR & \to & [0, 1]              \\
            t    & \tp & \bbP(X \leqslant t)
        \end{array}
    \]
\end{deff}

\begin{prop}
    Toute fonction de répartition $F_X$ vérifie
    \begin{enumerate}
        \item $F_X$ est croissante. ($\forall x \leqslant y, F_X(x) \leqslant F_X(y)$)
        \item $F_X$ est continue à droite. ($\forall t \in \bbR, \lim\limits_{s \to t^+} F_X(s) = F_X(t)$)
        \item $\lim\limits_{x \to +\infty} F_X(x) = 1$ et $\lim\limits_{x \to -\infty} F_X(x) = 0$.
    \end{enumerate}
\end{prop}

\begin{proof}
    $\ $
    \begin{enumerate}
        \item Pour $x \leqslant y$, $\{ X \leqslant x \} \subseteq \{ X \leqslant y \}$,
              donc $\bbP(X \leqslant x) \leqslant \bbP(X \leqslant y)$.
              i.e. $F_X(x) \leqslant F_X(y)$.
        \item Soit $(x_n)$ une suite décroissante telle que $x_n \underset{n \to +\infty}{\to} t$.
              Alors $\{ X \leqslant x_n \}$ est une suite décroissante d'événements.
              Par la continuité de la probabilité,
              \[ \bbP(X \leqslant t) = \lim_{n \to \infty} \downarrow \bbP(X \leqslant x_n)
                  = \bbP \left(\bigcap_{n \in \bbN} \{ X \leqslant x_n \}\right). \]
              Mais
              \[ \bigcap_{n \in \bbN} \{ X \leqslant x_n \} = \{ \forall n \in \bbN, X \leqslant x_n \}
                  = \{ X \leqslant \inf_{n \in \bbN} x_n \} = \{ X \leqslant t \}. \]
              Donc $F_X(t) = \lim\limits_{n \to \infty} \downarrow F_X(x_n)$.
        \item On remplace la suite $(x_n)$ dans la partie 2 par une suite $x_n \to +\infty$ croissante,
              ou une suite $x_n \to -\infty$ décroissante.
              \begin{itemize}
                  \item Si $(x_n)$ est croissante et $x_n \to \infty$,
                        \[ \lim_{n \to \infty} \uparrow \bbP(X \leqslant x_n) =
                            \bbP \left(\bigcup_{n \in \bbN} \{ X \leqslant x_n \}\right) = \bbP(\Omega) = 1. \]
                  \item Si $(x_n)$ est décroissante et $x_n \to -\infty$,
                        \[ \lim_{n \to \infty} \downarrow \bbP(X \leqslant x_n) =
                            \bbP \left(\bigcap_{n \in \bbN} \{ X \leqslant x_n \}\right) = \bbP(\varnothing) = 0. \qedhere \]
              \end{itemize}
    \end{enumerate}
\end{proof}

\begin{thm}[admis]
    La fonction de répartition d'une variable aléatoire réelle détermine sa loi.
\end{thm}

\begin{proof}[Idée de la preuve] $\ $
    \begin{itemize}
        \item $F_X$ détermine $\bbP(X \in ]-\infty, t])$.
        \item $B(\bbR) = \sigma(\{ ]-\infty, t]: t \in \bbR \})$. \qedhere
    \end{itemize}
\end{proof}

\begin{prop}
    Soit $X$ une variable aléatoire.
    Alors pour tout $x \in \bbR$,
    \[ \bbP(X = x) = F_X(x) - \lim_{y \to x^-} F_X(y). \]
\end{prop}

\begin{proof}
    \[ \{ X = x \} = \{ X \leqslant x \} \setminus \{ X < x \} =
        \{ X \leqslant x \} \setminus \left( \bigcup_{n \in \bbN} \{ X \leqslant x - \tfrac{1}{n} \} \right). \]
    Donc
    \begin{align*}
        \bbP(X = x) & = \bbP(X \leqslant x) - \lim_{n \to \infty} \uparrow \bbP\left(X \leqslant x - \tfrac{1}{n}\right) \\
                    & = F_X(x) - \lim_{n \to \infty} \uparrow F_X\left(x - \tfrac{1}{n}\right)                           \\
                    & = F_X(x) - \lim_{y \to x^-} F_X(y). \qedhere
    \end{align*}
\end{proof}

\begin{rqq}[Interprétation]
    Pour tout $x \in \bbR$,
    la loi $\bbP_X$ donne une mesure non nulle au singleton $\{ x \}$ ssi $F_X$ a une discontinuité à $x$.
\end{rqq}

\begin{rqq}
    Si $\bbP_X$ donne une mesure nulle à un intervalle $]a, b[$,
                    alors $F_X$ est constante sur $]a, b[$.
                    Pour tout $x < y$ et $x, y \in ]a, b[$,
    \[ F_X(y) - F_X(x) = \bbP(x < X \leqslant y) \leqslant \bbP(X \in ]a, b[) = \bbP_X(]a, b[) = 0. \]
\end{rqq}

\subsection{Fonction de répartition des variables aléatoires discrètes}

\begin{prop}
    Soient $X, Y$ deux variables aléatoires à valeurs dans $(E, \calE)$.
    Soit $f: (E, \calE) \to (F, \calF)$ une fonction mesurable.
    Si $X \overset{\text{loi}}{=} Y$, alors $f(X) \overset{\text{loi}}{=} f(Y)$.
\end{prop}

\begin{rqq}
    Autrement dit, la loi de $f(X)$ ne dépend que de la loi de $X$.
\end{rqq}

\begin{proof}
    Soit $C \in \calF$. On a $f(X) = f \circ X$.
    Alors \[ \bbP_{f(X)}(C) = \bbP(f(X) \in C) = \bbP(\{ \omega \in \Omega \mid f(X(\omega)) \in C \})
        = \bbP(\{ \omega \in \Omega \mid X(\omega) \in f^{-1}(C) \})
        = \bbP_X(f^{-1}(C)). \qedhere \]
\end{proof}

\section{Lois discrètes usuelles}

\subsection{Loi uniforme discrète}

Soit $E$ un ensemble non vide.

\begin{deff}
    Une variable aléatoire $X$ à valeurs dans $(E, \calP(E))$ suit la loi uniforme sur $E$,
    si $\forall x \in E, \bbP(X = x) = \frac{1}{\lrvert{E}}$.
    Alors on note $X \sim \unif(E)$. (``$X$ suit la loi uniforme sur $E$.'')
\end{deff}

\begin{exem}
    Si $a, b \in \bbZ$ avec $a < b$ et si $X \sim \unif(\llb a, b \rrb)$,
    alors $\forall a \leqslant k \leqslant b, \bbP(X = k) = \frac{1}{\lrvert{b - a + 1}}$.
\end{exem}

\subsection{Loi de Bernoulli}

Soit $p \in [0, 1]$ (paramètre).

\begin{deff}
    Une variable aléatoire à valeurs dans $\{ 0, 1 \}$ suit la loi de Bernoulli de paramètre $p$,
    si $\bbP(X = 1) = p, \bbP(X = 0) = 1 - p$.
    Alors on note $X \sim \ber(p)$. (``$X$ suit la loi de Bernoulli de paramètre $p$.'')
\end{deff}

\begin{exem}
    Si $A$ est un événement de probabilité $p$,
    alors $\bbI_A \sim \ber(p)$.
\end{exem}

\subsection{Loi binomiale}

Soient $n \in \bbN^*$ et $p \in [0, 1]$.

\begin{deff}
    Une variable aléatoire à valeurs dans $\llb 0, n \rrb$ suit la loi binomiale de paramètres $n$ et $p$,
    si $\forall k \in \llb 0, n \rrb, \bbP(X = k) = \binom{n}{k} p^k(1 - p)^{n - k}$.
    Alors on note $X \sim \calB(n, p)$.
\end{deff}

\begin{rqq}
    Cette loi est bien définie, car
    \[ \sum_{k = 0}^n \bbP(X = k) = \sum_{k = 0}^n \binom{n}{k} p^k(1 - p)^{n - k} = 1. \]
\end{rqq}

\begin{rqq}[Interprétation]
    La loi binomiale est la loi du nombre de ``succès'' dans $n$ épreuves / tests de Bernoulli indépendantes
    qui ont la même probabilité de succès $p$.
    Càd~: supposons $A_1, A_2, \ldots, A_n$~: $n$ événements mutuellement indépendants tels que
    $\forall k \in \llb 1, n \rrb, \bbP(A_k) = p$, et
    \[ X = \card \{ k \in \llb 1, n \rrb \mid A_k \text{ est réalisé} \} = \sum_{k = 1}^n \bbI_{A_k}. \]
    Alors $X \sim \calB(n, p)$.
\end{rqq}

\subsection{Loi géométrique}
Soit $p \in ]0, 1[$.

\begin{deff}
    Une variable aléatoire à valeurs dans $\bbN^*$ suit la loi géométrique de paramètre $p$,
    si $\forall k \in \bbN^*, \bbP(X = k) = p(1 - p)^{k - 1}$.
    Alors on note $X \sim \calG(p)$.
\end{deff}

\begin{rqq}[Interprétation]
    $X$ représente le rang (l'indice) du premier succès dans une suite (infinie)
    d'épreuves de Bernoulli indépendantes de même probabilité de succès $p$.
    Càd si
    \begin{itemize}
        \item $A_1, A_2, \ldots$ suite d'événements indépendants.
        \item $\forall n \in \bbN^*, \bbP(A_n) = p$.
        \item $X = \inf \{ n \in \bbN^* \mid A_n \text{ est réalisée} \}$.
    \end{itemize}
    Alors $X \sim \calG(p)$.
\end{rqq}

\begin{prop}[Absence de mémoire de la loi géométrique]
    Si $X \sim \calG(p)$, alors pour tout $k, l \in \bbN$,
    \begin{itemize}
        \item $\bbP(X > k) = (1 - p)^k$.
        \item $\bbP(X > k + l \mid X > k) = \bbP(X > l)$.
    \end{itemize}
\end{prop}

\begin{rqq}[Interprétation]
    \begin{itemize}
        \item $\{ X > k \}$~: le premier succès arrive après le temps $k$.
        \item $\{ X > k + l \}$~: le premier succès arrive après le temps $k + l$.
    \end{itemize}
\end{rqq}

\begin{proof}
    $\ $
    \begin{itemize}
        \item \[ \bbP(X > k) = \bbP \left( \bigcup_{n > k} \{ X = n \} \right)
                  = \sum_{n = k + 1}^\infty \bbP(X = n) = \sum_{n = k + 1}^\infty p(1 - p)^{n - 1} = (1 - p)^k \]
        \item \[ \bbP(X > k + l \mid X > k) = \frac{\bbP(X > k + l)}{\bbP(X > l)}
                  = \frac{(1 - p)^{k + l}}{(1 - p)^k} = (1 - p)^l = \bbP(X > l) \qedhere \]
    \end{itemize}
\end{proof}

\subsection{Loi de Poisson}

Soit $\lambda > 0$.

\begin{deff}
    Une variable aléatoire à valeurs dans $\bbN$ suit la loi de Poisson de paramètre $\lambda$,
    si $\forall k \in \bbN, \bbP(X = k) = e^{-\lambda} \frac{\lambda^k}{k!}$.
    Alors on note $X \sim \calP(\lambda)$.
\end{deff}

\begin{rqq}
    Cette loi est bien définie, car
    \[ \sum_{k = 0}^\infty e^{-\lambda} \frac{\lambda^k}{k!} = 1. \]
\end{rqq}

\begin{rqq}[Approximation de Poisson]
    Une loi binomiale comportant un grand nombre de tests,
    et une faible probabilité de succès est approximativement une loi de Poisson.
\end{rqq}

\begin{thm}
    Soit $(p_n)_{n \geqslant 1}$ une suite telle que
    \[ \forall n \in \bbN^*, p_n \in [0, 1], n \cdot p_n \underset{n \to \infty}{\to} \lambda > 0 \]

    Soit $X_n$ une variable aléatoire avec $X_n \sim \calB(n, p_n)$,
    alors $\forall k \in \bbN, \bbP(X_n = k) \underset{n \to \infty}{\to} e^{-\lambda} \frac{\lambda^k}{k!}$
\end{thm}

\begin{proof}
    Pour tout $n > k$, on a
    \[ \bbP(X_n = k) = \binom{n}{k} p_n^k(1 - p_n)^{n - k}
        = \frac{n(n - 1) \cdots (n - k + 1)}{k!} p_n^k(1 - p_n)^{n - k}. \]
    On peut réécrire
    \[ \bbP(X_n = k) = \frac{1}{k!} \cdot \frac{(np_n)^k}{(1 - p_n)^{k}} \cdot (1 - p_n)^n \cdot \frac{n(n - 1) \cdots (n - k + 1)}{n^k}. \]

    \begin{itemize}
        \item $\ln\left((1 - p_n)^n\right) = n \cdot \ln(1 - p_n) \underset{n \to \infty}{\sim} -n \cdot p_n
                  \underset{n \to \infty}{\to} -\lambda$, donc $(1 - p_n)^n \underset{n \to \infty}{\sim} e^{-\lambda}$.
        \item $\frac{n(n - 1) \cdots (n - k + 1)}{n^k} = \prod_{j = 0}^{k - 1} \frac{n - j}{n} \underset{n \to \infty}{\to} 1$.
    \end{itemize}

    On en déduit que $\bbP(X_n = k) \underset{n \to \infty}{\to} \frac{1}{k!} \lambda^k e^{-\lambda}$.
\end{proof}

\section{Vecteurs aléatoires et lois conjointes}

\begin{deff}
    Pour $n \geqslant 1$, on pose
    \[ B(\bbR^n) = \sigma(\{ I_1 \times I_2 \times \cdots \times I_n \mid
        I_1, I_2, \ldots, I_n \text{ sont des intervalles de } \bbR \}). \]

    On admet que $B(\bbR^n)$ contient tous les ensembles ouverts et tous les ensembles fermés de $\bbR^n$.
\end{deff}

$B(\bbR^n)$ est appelée la tribu borélienne de $\bbR^n$.

\begin{deff}
    Un \textbf{vecteur aléatoire} de dimension $n$ est une application mesurable
    \[ V: (\Omega, \calA, \bbP) \to (\bbR^n, B(\bbR^n)) \]
\end{deff}

\begin{thm}
    Soit $X_1, X_2, \ldots, X_n: \Omega \to \bbR$.
    L'application $X:
        \begin{array}[t]{ccc}
            \Omega & \to & \bbR^n                                                     \\
            \omega & \tp & \left(X_1(\omega), X_2(\omega), \ldots, X_n(\omega)\right)
        \end{array}
    $ est un vecteur aléatoire
    ssi $X_1, X_2, \ldots, X_n$ sont des variables aléatoires réelles.
\end{thm}

\begin{rqq}
    Autrement dit~:
    $X$ est mesurable par rapport à $B(\bbR^n)$ ssi les $X_i$ sont mesurables par rapport à $B(\bbR)$.
\end{rqq}

\begin{lem}
    Toute fonction continue $f: \bbR^n \to \bbR^m$ est mesurable par rapport aux $B(\bbR^n)$ et $B(\bbR^m)$.
\end{lem}

\begin{proof}
    Par définition de la continuité,
    \[ \forall U \text{ ouvert dans } \bbR^m, f^{-1}(U) \text{ est ouvert dans } \bbR^n. \]

    Donc $f^{-1}(U) \in B(\bbR^n)$,
    car $B(\bbR^n) = \sigma(\{ U: \text{ ouvert de } \bbR^n \})$.
    Et on a
    \begin{align*}
        f^{-1}\left(\bigcap_{n \in \bbN} A_n\right) & = \bigcap_{n \in \bbN} f^{-1}(A_n), \\
        f^{-1}\left(\bigcup_{n \in \bbN} A_n\right) & = \bigcup_{n \in \bbN} f^{-1}(A_n), \\
        f^{-1}\left(A^C\right)                      & = \left(f^{-1}(A)\right)^C.
    \end{align*}

    Donc $\{ A \subseteq \bbR^m \mid f^{-1}(A) \in B(\bbR^n) \}$ est une tribu,
    Donc $B(\bbR^m) \subseteq \{ A \subseteq \bbR^m \mid f^{-1}(A) \in B(\bbR^n) \}$.
\end{proof}

\begin{proof}[Démonstration du théorème]
    ``$\Rightarrow$'':
    Pour tout $k \in \llb 1, n \rrb, \pi_k:
        \begin{array}[t]{ccc}
            \bbR^n                  & \to & \bbR \\
            (x_1, x_2, \ldots, x_n) & \tp & x_k
        \end{array}
    $ est une fonction continue, donc mesurable.
    Si $X: \Omega \to \bbR^n$ est mesurable (i.e. $X$ est un vecteur aléatoire),
    alors $\forall k \in \bbN^*, \pi_k \circ X: \Omega \to \bbR$ est mesurable,
    i.e. $X_k$ est une variable aléatoire réelle.

    ``$\Leftarrow$'':
    Supposons que $X_1, X_2, \ldots, X_n: \Omega \to \bbR$ sont mesurables.
    Soit $\calC = \{ B \subseteq \bbR^n \mid X^{-1}(B) \in \calA \}$,
    alors $\calC$ est une tribu.
    Donc pour montrer que $\calC \supseteq B(\bbR^n)$,
    il suffit de montrer que pour tous les intervalles $I_1, I_2, \ldots, I_n \subseteq \bbR$,
    on a $I_1 \times I_2 \times \cdots \times I_n \in \calC$.
    Mais \[ X^{-1}(I_1 \times I_2 \times \cdots \times I_n)
        = \{ \omega \in \Omega \mid \forall k \in \llb 1, n \rrb, X_k(\omega) \in I_k \}
        = \bigcap\limits_{k = 1}^n \{ X_k \in I_k \}
        = \bigcap\limits_{k = 1}^n X_k^{-1}(I_k) \in \calA \]
    Donc $I_1 \times I_2 \times \cdots \times I_n \in \calC$.
    Puisque $B(\bbR^n) = \sigma(\{ I_1 \times I_2 \times \cdots \times I_n \mid
        I_1, I_2, \ldots, I_n \text{ sont des intervalles} \})$,
    on en déduit $\calC \supseteq B(\bbR^n)$.
    Donc $X: (\Omega, \calA) \to (\bbR^n, B(\bbR^n))$ est mesurable. \qedhere
\end{proof}

\begin{cor}
    Si $X_1, X_2, \ldots, X_n$ sont des variables aléatoires réelles et $f: \bbR^n \to \bbR^m$ est continue,
    alors $f(X_1, X_2, \ldots, X_n):
        \begin{array}[t]{ccc}
            \Omega & \to & \bbR^m                                           \\
            \omega & \tp & f(X_1(\omega), X_2(\omega), \ldots, X_n(\omega))
        \end{array}
    $ est un vecteur aléatoire.
\end{cor}

\begin{proof}
    Parce que la composée de deux fonctions mesurables est mesurable,
    \[ \Omega \overset{X = (X_1, X_2, \ldots, X_n)}{\to} \bbR^n \overset{f}{\to} \bbR^m \qedhere \]
\end{proof}

\begin{rqq}
    Si $X, Y$ sont deux variables aléatoires réelles,
    alors $X + Y, X - Y, XY, \lrvert{X}, \min(X, Y), \max(X, Y)$ sont des variables aléatoires réelles.
\end{rqq}

\begin{deff}
    Deux variables aléatoires $X, Y$ définies sur le même espace de probabilité sont dites égales presque sûrement,
    si $\bbP(X = Y) = 1$.
\end{deff}

\begin{rqq}
    Càd $X, Y: (\Omega, \calA, \bbP) \to E$,
    on dit $X = Y$ presque sûrement si $\bbP(\{ \omega \mid X(\omega) = Y(\omega) \}) = 1$.
\end{rqq}

\begin{prop}
    Si $X = Y$ presque sûrement, alors $X \overset{\text{loi}}{=} Y$.
\end{prop}

\begin{proof}
    Soit $N = \{ X \neq Y \} = \{ \omega \in \Omega \mid X(\omega) \neq Y(\omega) \}$.
    Soit $B \in \calE$.

    Alors $\bbP_X(B) = \bbP(\{ X \in B \}) = \bbP(\{ X \in B \} \setminus N) + \bbP(\{ X \in B \} \cap N)$.

    D'une part,
    \[ \{ X \in B \} \setminus N = \{ X \in B \} \cap \{ X = Y \}
        = \{ Y \in B \} \cap \{ X = Y \} = \{ Y \in B \} \setminus N, \]
    donc $\bbP(\{ X \in B \} \setminus N) = \bbP(\{ Y \in B \} \setminus N)$.

    D'autre part, $\{ X \in B \} \cap N \subseteq N$,
    et $\bbP(N) = \bbP(X \neq Y) = 0$.

    Donc
    \[ \bbP_X(B) = \bbP(\{ X \in B \} \setminus N) = \bbP(\{ Y \in B \} \setminus N) = \bbP_Y(B). \]

    Donc $\bbP_X = \bbP_Y$, càd $X \overset{\text{loi}}{=} Y$. \qedhere
\end{proof}

\begin{rqq}
    Problème~: en général, $\{ X = Y \}$ n'est pas nécessairement mesurable.

    \begin{itemize}
        \item Cas réel~: si $X, Y$ sont à valeurs dans $(\bbR^n, B(\bbR^n))$,
              alors $\{ X = Y \} = \{ X - Y \in \{0_{\bbR^n}\} \}$.
        \item Cas général~: on définit $X = Y$ presque sûrement ssi
              \[ \exists A \in \calA, A \subseteq \{ X = Y \}, \bbP(A) = 1. \]
    \end{itemize}
\end{rqq}

\begin{deff}
    La \textbf{loi conjointe} des variables aléatoires réelles $X_1, X_2, \ldots, X_n$
    est la loi du vecteur aléatoire $(X_1, X_2, \ldots, X_n)$.
    Elle est notée $\bbP_{(X_1, X_2, \ldots, X_n)}$ (mesure de probabilité sur $\bbR^n$).
    Dans ce contexte, les lois de $X_1, X_2, \ldots, X_n$ (i.e. $\bbP_{X_1}, \bbP_{X_2}, \ldots, \bbP_{X_n}$)
    sont appelées les lois marginales du vecteur aléatoire.
\end{deff}

\begin{exer}
    Pour $A, B \in B(\bbR)$, et $X, Y$ deux variables aléatoires réelles,
    on a
    \begin{align*}
        \bbP_{(X, Y)} (A \times B) & = \bbP(X \in A \text{ et } Y \in B), \\
        \bbP_X(A)                  & = \bbP_{(X, Y)}(A \times \bbR),      \\
        \bbP_Y(B)                  & = \bbP_{(X, Y)}(\bbR \times B).
    \end{align*}
\end{exer}

\begin{rqq}[Cas discret]
    Soient $X, Y$ deux variables discrètes.
    Soient $D = X(\Omega), D' = Y(\Omega)$ ($D, D'$ sont au plus dénombrables).
    Soit $p_{xy} = \bbP(X = x, Y = y), \forall (x, y) \in D \times D'$.
    Alors on a
    \begin{align*}
         & \sum_{(x, y) \in D \times D'} p_{xy} = 1,                    \\
         & \forall x \in D, \quad \bbP(X = x) = \sum_{y \in D'} p_{xy}, \\
         & \forall y \in D', \quad \bbP(Y = y) = \sum_{x \in D} p_{xy}.
    \end{align*}
\end{rqq}

\begin{proof}
    \[ \bbP(Y = y) = \bbP \left(\bigsqcup_{x \in D} \{ X = x, Y = y \}\right) = \sum_{x \in D} \bbP(X = x, Y = y). \qedhere \]
\end{proof}

\subsection{Loi conditionnelle sachant une valeur de probabilité positive}

Soient $X, Y$ deux variables aléatoires, et $y$ une valeur de $Y$ telle que $\bbP(Y = y) > 0$.

\begin{deff}
    La loi conditionnelle de $X$ sachant $Y = y$ est définie par
    \[ \forall A \in \calE, \bbP_{X \mid Y = y}(A) = \bbP(X \in A \mid Y = y) \]
\end{deff}

\begin{rqq}
    Càd
    \[ \bbP_{X \mid Y = y}(A) = \frac{\bbP(\{ X \in A \} \cap \{ Y = y \})}{\bbP(Y = y)} \]
\end{rqq}

\begin{exem}
    On lance deux dés équilibrés.
    Soient $X = \text{résultat du premier dé} \in \llb 1, 6 \rrb$,
    $S = \text{la somme des résultats des deux dés} \in \llb 2, 12 \rrb$.
    On considère l'événement $\{ S = 7 \}$.

    Pour tout $k \in \llb 1, 6 \rrb$,
    \[ \bbP(X = k, S = 7) = \bbP(X = k) \cdot \bbP(S - X = 7 - k) = \left(\frac{1}{6}\right)^2. \]
    Donc $\bbP(S = 7) = \sum\limits_{k = 1}^6 \bbP(X = k, S = 7) = \frac{1}{6}$, et
    \[ \bbP(X = k \mid S = 7) = \frac{\frac{1}{36}}{\frac{1}{6}} = \frac{1}{6}. \]
    Càd la loi conditionnelle $\bbP_{X \mid S = 7}$ est la même que la loi de $X$,
    càd la loi uniforme sur $\llb 1, 6 \rrb$~:
    \[ \bbP_{X \mid S = 7} = \bbP_X = \unif(\llb 1, 6 \rrb). \]
\end{exem}

\section{Indépendance de variables aléatoires}

Soit $(\Omega, \calF, \bbP)$ un espace de probabilité.

\begin{deff}[Indépendance de deux événements] $\ $
    \begin{itemize}
        \item On dit que $A, B$ sont indépendants, si
        \[ \bbP(A \cap B) = \bbP(A) \bbP(B). \]
        \item On dit que $n$ événements $A_1, A_2, \ldots, A_n$ sont mutuellement indépendants, si
        \[ \forall I \subseteq \llb 1, n \rrb, \bbP \left(\bigcap_{i \in I} A_i\right) = \prod_{i \in I} \bbP(A_i). \]
    \end{itemize}
\end{deff}

\begin{deff}[Indépendance de deux variables aléatoires] $\ $
    Soit $X: (\Omega, \calF) \to (\bbR^n, \calB(\bbR^n))$
    et $Y: (\Omega, \calF) \to (\bbR^m, \calB(\bbR^m))$.

    On dit que $X$ et $Y$ sont indépendantes si pour tous $A \in \calB(\bbR^n), B \in \calB(\bbR^m)$,
    \[ \bbP(X \in A, Y \in B) = \bbP(X \in A) \bbP(Y \in B). \]

    On note $X \perp Y$.
\end{deff}

\begin{rqq}[Conséquence]
    Si $X \perp Y$, alors $\bbP(X \in A \mid Y \in B) = \bbP(X \in A), (\bbP(Y \in B) > 0)$.
\end{rqq}

\begin{prop}
    Si $X, Y$ sont discrètes, alors $X \perp Y$ ssi
    \[ \forall x \in X(\Omega), y \in Y(\Omega), \bbP(X = x, Y = y) = \bbP(X = x) \bbP(Y = y) \]
\end{prop}

\begin{proof}
    Le sens ``$\Rightarrow$'' s'obtient en prenant
    \[ A = \{ x \}, B = \{ y \} \]

    ``$\Leftarrow$''~: Réciproquement, posons $D = X(\Omega), D' = Y(\Omega)$, $D, D'$ sont au plus dénombrables.
    Pour tous ensembles mesurables $A, B$, on a
    \begin{align*}
        \bbP(X \in A, Y \in B) &= \sum_{\substack{x \in D \\ y \in D'}} \bbP(X = x, Y = y)
        = \sum_{\substack{x \in D \\ y \in D'}} \bbP(X = x) \bbP(Y = y) \\
        &= \left(\sum_{x \in D} \bbP(X = x)\right)\left(\sum_{y \in D'} \bbP(Y = y)\right)
        = \bbP(X \in A) \bbP(Y \in B).
    \end{align*}
\end{proof}

\begin{prop}[Indépendance par les fonctions de répartition]
    Deux variables $X, Y$ (à valeurs réelles) sont indépendantes ssi
    \[ F_{X, Y}(x, y) = F_X(x) F_Y(y), \quad \forall x, y \in \bbR, \]
    c'est-à-dire
    \[ \bbP(X \leqslant x, Y \leqslant y) = \bbP(X \leqslant x) \bbP(Y \leqslant y). \]
\end{prop}

\begin{prop}[Indépendance par transformation]
    Soient $X, Y$ deux variables aléatoires indépendantes.
    Soit $f, g: \bbR^n \to \bbR$ deux applications mesurables.
    Alors $f(X)$ et $g(Y)$ sont indépendantes.
\end{prop}

\begin{proof}
    Pour tous ensembles mesurables $A, B \in \calB(\bbR^n)$,
    \[ \bbP(f(X) \in A, g(Y) \in B) = \bbP(X \in f^{-1}(A), Y \in g^{-1}(B))
    = \bbP(X \in f^{-1}(A)) \bbP(Y \in g^{-1}(B)) = \bbP(f(X) \in A) \bbP(g(Y) \in B). \]

    Donc $f(X)$ et $g(Y)$ sont indépendantes.
\end{proof}

\subsection{Indépendance d'une famille de variables}

Soit $I$ un ensemble d'indices.

\begin{deff}
    Soit $(X_i)_{i \in I}$ une famille de variables aléatoires.
    On dit que $(X_i)_{i \in I}$ sont mutuellement indépendantes,
    si pour toute partie finie $J \subseteq I$,
    pour tous ensembles mesurables $A_j$, on a
    \[ \bbP \left(\bigcap_{j \in J} A_j\right) = \prod_{j \in J} \bbP(A_j) \]
\end{deff}

\begin{rqq}
    $(X_i)_{i \in I}$ sont indépendantes deux à deux,
    si pour tout $i \neq j$, $X_i$ et $X_j$ sont indépendantes.
\end{rqq}

\begin{exem}
    Soit $\Omega = \{ 0, 1 \}$.

    \begin{table}[htb]
        \centering
        \begin{tabular}{c|c|c}
            \diagbox{X}{Y} & 0           & 1           \\
            \hline
            0   & $\frac{1}{4}$ & $\frac{1}{4}$ \\
            \hline
            1   & $\frac{1}{4}$ & $\frac{1}{4}$
        \end{tabular}
    \end{table}

    \[ \bbP(X = i, Y = j) = \frac{1}{2^2}, \quad i, j \in \{ 0, 1 \}. \]
    Les lois marginales sont
    \[ \bbP(X = 0) = \bbP(X = 1) = \bbP(Y = 0) = \bbP(Y = 1) = \frac{1}{2}. \]
    
    Donc $X \perp Y$.
    
    On définit $Z = \bbI_{X = Y}$, alors
    \[ \bbP(Z = 1) = \bbP(Z = 0) = \frac{1}{2}. \]

    Et
    \begin{align*}
        \bbP(X = 0, Z = 0) &= \bbP(X = 0, Y = 1) = \frac{1}{4}, &
        \bbP(X = 0, Z = 1) &= \bbP(X = 0, Y = 0) = \frac{1}{4}, \\
        \bbP(X = 1, Z = 0) &= \bbP(X = 1, Y = 0) = \frac{1}{4}, &
        \bbP(X = 1, Z = 1) &= \bbP(X = 1, Y = 1) = \frac{1}{4}.
    \end{align*}

    Donc $X \perp Z$, par symétrie, $Y \perp Z$.
    Mais $X, Y, Z$ ne sont pas mutuellement indépendantes.
\end{exem}

\begin{prop}
    Soient $X_1, X_2, \ldots, X_n$ variables aléatoires réelles mutuellement indépendantes.
    Si $f: \bbR^p \to \bbR$, et $g: \bbR^{n - p} \to \bbR$ sont mesurables,
    alors les deux variables aléatoires
    \[ f(X_1, X_2, \ldots, X_p), g(X_{p + 1}, X_{p + 2}, \ldots, X_n) \]
    sont indépendantes.
\end{prop}

\begin{proof}
    Posons $U = (X_1, X_2, \ldots, X_p) \in \bbR^p, V = (X_{p + 1}, X_{p + 2}, \ldots, X_n) \in \bbR^{n - p}$.
    On va démontrer que $U \perp V$.
    On prend $A = I_1 \times I_2 \times \cdots \times I_p, B = I_{p + 1} \times I_{p + 2} \times \cdots \times I_n$,
    où $I_j$ sont des intervalles dans $\bbR$.
    Par l'indépendance,
    \begin{align*}
        \bbP(U \in A, V \in B) &= \bbP(X_1 \in I_1, \ldots, X_p \in I_p, X_{p + 1} \in I_{p + 1}, \ldots, X_n \in I_n) \\
        &= \bbP(X_1 \in I_1, \ldots, X_p \in I_p) \cdot \bbP(X_{p + 1} \in I_{p + 1}, \ldots, X_n \in I_n) \\
        &= \bbP(U \in A) \bbP(V \in B).
    \end{align*}

    Comme les produits d'intervalles engendrent les tribus boréliennes,
    cette égalité se généralise à tous $A \in \calB(\bbR^p), B \in \calB(\bbR^{n - p})$.
    Donc $U \perp V$.
\end{proof}

Dans la suite de cours, nous admettons le théorème suivant.
\begin{thm}[Existence d'une suite de variables aléatoires indépendantes]
    Pour toute suite $(P_n)$ de lois de probabilité,
    il existe un espace de probabilité $(\Omega, \calF, \bbP)$,
    et une suite $(X_n)$ de variables aléatoires indépendantes sur $(\Omega, \calF, \bbP)$,
    telle que pour tout $n$, la loi de $X_n$ soit $P_n$.
\end{thm}

\begin{rqq}
    En pratique, si $P_n = P$, pour tout $n$,
    (on dit que $(X_n)$ est une suite de variables aléatoires
    indépendantes, identiquement distribuées ($(X_n)$ i.i.d.)),
    une telle suite modélise une suite d'épreuves identiques aux résultats indépendants.
\end{rqq}

\subsection{Applications}

\paragraph{Convolution dans le cas discret}

Soient $f, g: \bbZ \to \bbR$. La convolution discrète est définie par
\[ (f * g)(n) = \sum_{k \in \bbZ} f(n - k) g(k). \]

convolution discrète pour $f, g$ telles que
\[ \sum\limits_{n} \lrvert{f(n)} < +\infty, \sum\limits_{n} \lrvert{g(n)} < +\infty \]

Soient $X, Y$ deux variables aléatoires indépendantes à valeurs dans $\bbZ$,
avec $\mathrm{Loi}(X) = {}p$ et $\mathrm{Loi}(Y) = {}q$, c'est-à-dire
$\bbP(X = n) = {}p(n)$ et $\bbP(Y = n) = {}q(n)$, où
${}p, {}q: \bbZ \to [0, +\infty[$ satisfont $\sum_{n \in \bbZ} {}p(n) = \sum_{n \in \bbZ} {}q(n) = 1$.

Question: $\mathrm{Loi}(X + Y) = ?$

On a
\[ \bbP(X + Y = n) = \sum_{k \in \bbZ} \bbP(X = n - k, Y = k)
= \sum_{k \in \bbZ} {}p(n - k) {}q(k) = ({}p * {}q)(n). \]
Donc $\mathrm{Loi}(X + Y) = {}p * {}q$, la convolution discrète de $\mathrm{Loi}(X)$ et $\mathrm{Loi}(Y)$.

\paragraph{Loi binomiale}

\begin{prop}
    Soient $X_1, X_2, \ldots, X_n$ des variables indépendantes telles que $X_i \sim \ber(p)$ avec $p \in ]0, 1[$.
    Alors $Y = \sum_{k = 1}^n X_k \sim \calB(n, p)$.
    Plus précisément, $\bbP(Y = k) = \binom{n}{k} p^k (1-p)^{n-k}, \forall k \in \llb 0, n \rrb$.
\end{prop}

\begin{proof}
    $\bbP(X_i = 1) = p, \bbP(X_i = 0) = 1 - p$.
    La somme $\sum_{k = 1}^n X_k$ compte le nombre de succès,
    pour $k \in \llb 1, n \rrb, \{ Y = k \} = \{ X_1 + X_2 + \cdots + X_n = k \}$
    est l'union disjointe sur les parties $J \subseteq \llb 1, n \rrb$ de cardinal $k$
    des événements $\left(\bigcap_{i \in J} \{ X_i = 1 \}\right) \cap \left(\bigcap_{i \notin J} \{ X_i = 0 \}\right)$
    pour chaque $J$, on a $\# J = k, \bbP((\bigcap\limits_{i \in J} \{ X_i = 1 \}) \cap (\bigcap\limits_{i \notin J} \{ X_i = 0 \})) = p^k q^{n-k}$
    Donc $\bbP(Y = k) = \binom{n}{k} p^k q^{n - k} = \binom{n}{k} p^k (1-p)^{n-k}$.
\end{proof}

\begin{proof}
    Raisonnement par récurrence.

    Il s'agit d'une propriété $\calP(n)$, on procède en deux étapes.

    Étape 1 (Initialisation)~: $\calP(n)$ est vraie pour $n = n_0$.

    Étape 2 (Hérédité)~: sachant que $\calP(n)$ est vraie, on en déduit que $\calP(n + 1)$ est vraie.

    Conclusion: $\calP(n)$ est vraie pour tout $n$ entier, $n \geqslant n_0$.

    Pour notre exemple.
    \[ \calP(n): \bbP(Y_n = k) = \binom{n}{k}, \forall k \in \llb 0, n \rrb \]

    Initialisation: $n = 1, k \in \{ 0, 1 \}$.
    \[ \bbP(Y_1 = 0) = \bbP(X_1 = 0) = q, \bbP(Y_1 = 1) = \bbP(X_1 = 1) = p \]
    Donc $\calP(1)$ est vraie.

    Hérédité~: supposons que $\calP(n)$ est vraie,
    \[ Y_{n + 1} = Y_n + X_{n + 1} \]
    où $Y_n$ et $X_{n + 1}$ sont indépendantes.
    Pour $1 \leqslant k \leqslant n + 1$, on a
    \begin{align*}
        \bbP(Y_{n + 1} = k) &= \bbP(Y_n = k - 1) \bbP(X_{n + 1} = 1) + \bbP(Y_n = k) \bbP(X_{n + 1} = 0) \\
        &= \binom{n}{k - 1} p^{k - 1} q^{n - k + 1} p + \binom{n}{k} p^k q^{n - k} q \\
        &= \left(\binom{n}{k - 1} + \binom{n}{k}\right) p^k q^{n - k} \\
        &= \binom{n + 1}{k} p^k q^{n - k}.
    \end{align*}
    
    Et
    \[ \bbP(Y_{n + 1} = 0) = \bbP(X_1 = 0, X_2 = 0, \ldots, X_{n + 1} = 0) = q^{n + 1} \]

    Par conséquent, $\calP(n + 1)$ est vraie.

    En conclusion, $\calP(n)$ est vraie pour tout $n \in \bbN^*$.
\end{proof}